% 原书第 613 页（PDF 第 626 页）
\BookPage{613}{626}
\BookHeading{读文渊阁《四库全书》本《明史‧土司传考证》}{读文渊阁《四库全书》本《明史‧土司传考证》}{戴晋新*}
乾隆四年（1739）刊行的清代官修《明史》与乾隆五十四年（1789）收入《四库全书》的《明史》是两个不同的《明史》版本，前者刊刻于武英殿，流传甚广，为一般通行的版本，1974年中华书局刊行的标点校刊本（以下称标校本）《明史》依据的即为这个版本；后者在清代被称为「新定之本」1，后来通称「四库本」。《四库全书》抄本分藏七处，本文讨论的是文渊阁四库本《明史》。两个《明史》版本不仅刊行的时间不同，内容上也有一些明显的差异。乾隆朝纂辑《四库全书》，所收各书例经审校，《明史》亦无例外。乾隆四十年（1775）五月，清高宗谕令改纂四库拟收的《明纪纲目》2；同月，高宗又下令将《明史》一并查考改正3；乾隆四十二（1777)年五月，高宗更下令将《明史》本纪逐一考核添修，「候朕亲阅审定」4。在皇帝的命令下，四库馆臣对《明史》进行了查考改正的工作，实际的成果反映在本纪与列传上。检视文渊阁四库本《明史》与武英殿版《明史》，可以发现四库本《明史》本纪部分有增补史事、详其原委、润饰文字、修订赞语等情形；列传部分也有注释说明、考异存疑、辨误定讹、更正用语等情形，传文既有所修改，卷末并附有考证。论者因此认为四库本《明史》「是清朝官方纂修《明史》的最终成果」，卷末的考证内容

\vfill
\begin{BookNotes}
\BookNote{*}{作者任职于辅仁大学历史学系。}
\BookNote{1.}{《四库全书总目》、《清文献通考》俱如此称。}
\BookNote{2.}{《清高宗实录》卷九八二，乾隆四十年五月辛酉。}
\BookNote{3.}{《清高宗实录》卷九八三，乾隆四十年五月甲子。}
\BookNote{4.}{《清高宗实录》卷一三二，乾隆四十二年五月丁丑。}
\BookNote{5.}{参：乔治忠、杨艳秋，〈《四库全书》本《明史》发覆〉，《清史研究》1999年第4期，页68-70；又见:乔治忠，《中国官方史学与私家史学》（北京:北京图书馆出版社，2008)，页287-299。}
\end{BookNotes}

% 原书第 614 页（PDF 第 627 页）
\BookPage{614}{627}
\noindent 「是清朝官修《明史》最后定本的组成部分」。过去由于《四库全书》不易得见，这部内文经过修订并附有考证的四库本《明史》鲜见有人探用，即使1986年台湾商务印书馆将文渊阁《四库全书》影印出版后，情况似乎也没有明显的改变；学界有关《明史》的征引还是以武英殿版或中华书局的标校本为主，有关《明史》的考证往往仍以王颂蔚《明史考证捃逸》为前提。本文不拟复述《明史》版本的相关研究8，谨以文渊阁《四库全书》本《明史·土司传考证》为例，谈谈个人对它的认识。

二

四库本《明史》列传的卷末考证，每则皆录撰稿者姓名，如「臣某某按」。虽然光绪年间纂辑《明史考证捃逸》的王颂蔚说「进呈本所题姓名往往与正本不同」9，但每则考证的作者最后仍当以写定的四库本作为依据。四库本《明史·土司传考证》的撰稿人为严福与章宗瀛，严福只考了一卷3则，其余九卷150则皆由章宗瀛负责。严福（？一1792)，字景纯、井仁，江苏苏州府吴县人。乾隆四十年进士，乾隆四十三年起任翰林院编修至乾隆五十七年，历充四库全书校对官，武英殿、国史馆、方略馆纂修官，乾隆四十四年河南乡试正考官，乾隆五十七年上书房行走，同年卒。1°章宗瀛（生卒年不详），字登之，浙江绍兴府会稽县人。乾隆四十年进士，乾隆四十三年起任翰林院编修，历充乾隆四十五年顺天乡试同考官，乾隆四十六年会试同考官，乾隆四十八年顺天乡试同考官。乾隆五十年乾清宫考列三等，所有在四库全书馆议叙升用之处注销，仍罚俸半年。11「性耿介，和珅招之不屈，三十余年不迁官」12。两人

\vfill
\begin{BookNotes}
\BookNote{6.}{同上，页70；页299。}
\BookNote{7.}{《明史》考证不止一家，如黄云眉，《明史考证》（北京：中华书局，1979）亦颇负盛名，其他各种单篇、单卷的考证、笺注不胜枚举；但从成书背景与性质看，《明史》考证除四库本《明史考证》而外，自当以王书为基本前提。}
\BookNote{8.}{有关《明史》修撰经过与版本，论者甚多，最近的讨论可参：衣若兰，〈《明史》稿本探研—从万斯同《明史稿》到四库本《明史》〉，见:虞浩旭、饶国庆主编，《万斯同与明史:纪念万斯同诞辰370周年国际学术研讨会论文集（下），宁波：宁波出版社，2008。}
\BookNote{9.}{王颂蔚，《明史考证捃逸·各卷纂辑姓名》，引见《新校本明史并附编六种》第十二册（台北：鼎文书局，1975），明史附编四，页10。}
\BookNote{10.}{《〔同治〕苏州府志》卷八三，〈人物十〉。}
\BookNote{11.}{《清高宗实录》卷一二二四，乾隆五十年二月己丑。}
\BookNote{12.}{《〔民国〕绍兴县志资料第一辑·人物列传第一编》，页139。}
\end{BookNotes}

% 原书第 615 页（PDF 第 628 页）
\BookPage{615}{628}
\noindent 皆进士出身，四库馆臣，不过章宗瀛后来被除名，〈办理四库全书在事诸臣职名〉中并没有他的名字。

四库本《明史·土司传考证》分属10卷，合计153则，依其内容性质，可分为三类：

第一类考述：因传文未载、未晰、未详，考而述之、释之者，包括过程、原委、沿革，名词、谥号等，计88则。

第二类考订：因传文有讹误，考而定之者，包括时间、地名、人物人名、官名、隶属等，计36则。

第三类考异：因传文内容别有异说，存异待定者，包括时间、事实、名称、沿革等，计29则。

这些抄录在各卷卷末的考证文字，据王颂蔚的说法，当初有「进呈本」、「稿本」、「正本」、「初刊样本」13四种工作稿本。最后录进四库本《明史》的当为定稿，撰稿者署名，以示负责。经由数稿而成的定稿照说应有较高水平，然而仔细校读发现其中仍有许多失误，举例如次：

讹误:卷311松藩卫传十年黑虎塞贼首多儿太伏诛条，「塞」当为「寨」。卷313临

\begin{BookQuote}
安传曰思陀甸条，「地理志作思陀」，「陀」当为「佗」。卷315孟养传十四年复设孟养宣慰司条，「谨未载」当为「传未载」。卷316思南传元传宣慰司条，「元传」当为「元改」；「秦定中」当为「泰定中」。卷317平乐传府江有南岸三洞诸撞皆属荔浦延袤千余里条，「柳庆士民」当为「柳庆土民」。卷318上石西州传宋属太平条，「太平」当为「永平」。卷319向武州传左副将军魏观，「魏」当为「韩」。
\end{BookQuote}

失考：卷318镇安府传帝命两广总督陶谐治之条，传文「嘉靖十四年，……按臣曾守

\begin{BookQuote}
约以闻帝命新建伯王守仁治之」，《考证》指出王守仁卒于嘉靖八年（误，当为七年），不可能于嘉靖十四年受命；并据《明史·陶谐传》、《明实录》、《明史记事本末》，因「时陶谐总督两广」，遂改「新建伯王守仁」为「两广
\end{BookQuote}

\vfill
\begin{BookNotes}
\BookNote{13.}{王颂蔚，《明史考证捃逸·序》，引见《新校本明史并附编六种》第十二册（台北：鼎文书局，1975)，明史附编四，页1-2。}
\end{BookNotes}

% 原书第 616 页（PDF 第 629 页）
\BookPage{616}{629}
\begin{BookQuote}
总督陶谐」。事实上与本则考证相关的田州卢苏案，分别发生于嘉靖六年与嘉靖十四年，王守仁与嘉靖六年那次有关，见同卷思恩府传、田州传以及《世宗实录》嘉靖六年三月乙未条；镇安府传所叙为嘉靖十四年事，据《世宗实录》嘉靖十四年十二月丁未条谓「巡按御史曾守约以闻」，「命守臣亟为议处」，镇安府传「新建伯王守仁」当改为「守臣」，岂可臆改为「两广总督陶谐」？查〈陶谐传〉亦未载此事。《明史》本传误两事为一事，四库本《明史·土司传考证》并未发现，考证失据，一讹再讹。卷318田州传巡按御史曾守约以闻帝命守仁亟为勘处条，失考同前。
\end{BookQuote}

脱漏：卷311松藩卫传十年黑虎塞贼首多儿太伏诛条，「臣章宗瀛」下脱一「按」字。

衍误：卷311天全卫传可自碉门出柘木场条，「西口旧本误作自碉门出枯木任伤谨据

\begin{BookQuote}
改」，实系衍文。
\end{BookQuote}

笼统:卷317平乐传乃定善后六策以闻条，「善后六策明实录诸书均无考」，「诸书」未明何指。

失序:卷319思陵州传嘉靖十九年总督张经条，据卷205〈张经传〉改蔡经为张经；

\begin{BookQuote}
按卷318凭祥州传已见「总督蔡经」，四库本传文与《明史考证捃逸》皆改蔡经为张经；两传考证应详未详，应略未略，前后失序。又〈张经传〉谓「初冒姓蔡，久之乃复」，应一并考入，以明原委。
\end{BookQuote}

错乱:卷311宁番卫传二十七年改宁番卫条，据同卷建昌卫传、越巂卫传、明史地

\begin{BookQuote}
理志为证，未引统志，岂可谓「据统志改」？《明史考证捃逸》本条亦未引统志。卷311天全卫传可自碉门出柘木场条，「任伤柘木」，当为「枯木任伤」。卷312永宁卫传樊龙等以增行粮为名乘机反条、卷318田州府传猇亦自杀条、卷319龙州传都御史杨旦讨韦璋杀之条，各条所引「通鉴纪事」《明史考证捃逸》皆作「明史纪事本末」。卷314姚安传土舍凤文作乱条，应作「武定传土舍凤朝文作乱」，土司名与人名双误。
\end{BookQuote}

无征：卷314元江传条，未书所据。


% 原书第 617 页（PDF 第 630 页）
\BookPage{617}{630}
\noindent 不一:依例《明史》各卷但书卷名，卷311松藩卫传宣德七年升龙州为宣抚司条书

\begin{BookQuote}
「明史地理志」，应书「地理志」。
\end{BookQuote}

另外有些失误比较复杂，例如《明史·麓川传〉载：

\begin{BookQuote}
麓川、平缅，元时皆属缅甸。缅甸，古朱波地也。宋宁宗时，缅甸、波斯等国进白象，缅甸之名自此始。缅在云南之西南，最穷远。与八百国、占成接境。有城郭室屋，人皆楼居，地产象马。元时最强盛。元尝遣使招之，始入贡。
\end{BookQuote}

四库本《明史》卷314〈麓川传〉改正为：

\begin{BookQuote}
麓川、平缅，汉时为永昌郡，檄外。麓川之地曰大布茫，曰琰头附赛，曰琰中弹吉，曰琰尾福禄培，在茫施东；平缅之地曰骠琰，曰罗必四庄，曰小沙摩弄，曰骠琰头，在麓川东北；皆处金沙江之南，故为白蛮所据。其地又有龙川江，为麓川达木邦之路，蛮恃以为险。元中统初始内附，至元十三年立为路，并隶于金齿宣抚司。
\end{BookQuote}

这段文字见于《明史考证捃逸》，并有按语：

\begin{BookQuote}
按一统志及元史、经世大典、明史、明一统志、明史记事本末诸书，麓川、平缅在元时立两路，属金齿宣抚司，此作属缅甸，误，且以下接纪缅甸事实业已详载彼传，尤不应复叙。兹据诸书所载麓川、平缅将土建置改正。
\end{BookQuote}

《捃逸》所言甚是。此例《明史·麓川传》错的离谱，四库本大幅改写订正传文，却删去卷末考语，未予说明，做法失当。

三

除了上文所举失误，四库本《明史·土司传考证》的征引也值得观察。首先是征引文献名称体例不一与讹误的问题，四库本《明史·土司传考证》征引文献时，或仅录书名不录篇卷，如引《明实录》不录帝号，引《唐书》、《新唐书》、《宋史》不录卷名；或并录书名卷名，如称《汉书·地理志》、《元史·地理志》；书名或仅作省称，如《读史方舆纪要》称《方舆纪要》；同书省称不一者亦有之，如《明一统


% 原书第 618 页（PDF 第 631 页）
\BookPage{618}{631}
\noindent 志》省称《明统志》、《一统志》、《统志》；引《明史》依例但书卷名不称书名，如〈地理志〉、〈英宗本纪〉，唯亦有违例情形，如〈地理志〉书作《明史·地理志》，体例不一；书名亦有舛误，如《通鉴纪事》在《明史考证捃逸》中皆作《明史纪事本末》。

其次，征引资料种类与征引则数也透露出一些讯息，统计分析如下：

《明纪事本末》16则（即《明史纪事本末》）《蛮司合志》24则《通鉴纪事》4则（当为《明史纪事本末》）《峒溪纤志》1则《明实录》34则《实录》1则（当为《明实录》）《汉书·地理志》1则《元史·地理志》4则《新唐书》1则《唐书》2则《宋史》2则《明史·地理志》3则《统志》17则（当为《明一统志》）《明统志》3则（当为《明一统志》）《一统志》16则（当为《明一统志》）《方舆纪要》3则（即《读史方舆纪要》《识大录》5则《通志》1则（当为《明通志》）《明通志》4则《王守仁文集》1则《唐荆川集》1则《土夷考》4则


% 原书第 619 页（PDF 第 632 页）
\BookPage{619}{632}
\noindent 〈地理志〉24则〈英宗本纪〉1则〈杨畏知传〉1则<邹文盛传〉1则〈陶谐传〉1则〈秦纮传〉1则〈孟养土司本传〉1则〈方瑛传〉1则（卷316)〈田州传〉1则（卷318）<镇安传〉1则（卷318）<归顺州传〉1则（卷319）<沈希仪传〉1则（卷319）〈张经传〉1则（卷319）〈各传〉1则（卷319）

征引则数最多者，为《明实录》、《明一统志》、《明史·地理志》、《蛮司合志》、《明史纪事本末》、《明史》各传。若不论各书征引篇卷与则数，仅列书名，计得17种。同样不论征引篇卷与则数，仅列书名，中华书局标校本《明史·土司传·校勘记》征引材料计得18种。两相对照如次：

\BookListLine{四库全书本}{中华书局本}
\BookListLine{《汉书》}{《汉书》}
\BookListLine{《元史》}{《元史》}
\BookListLine{《唐书》}{《旧唐书》}
\BookListLine{《明实录》}{《明实录》}
\BookListLine{《明一统志》}{《明一统志》}
\BookListLine{《明史》}{《明史》}
\BookListLine{《读史方舆纪要》}{《读史方舆纪要》}
\BookListLine{《明通志》}{《寰宇通志》}
\BookListLine{《蛮司合志》}{《蛮司合志》}
\par


% 原书第 620 页（PDF 第 633 页）
\BookPage{620}{633}
\BookListLine{《明史纪事本末》}{《明史稿》}
\BookListLine{《峒溪纤志》}{《国榷》}
\BookListLine{《识大录》}{《明会典》}
\BookListLine{《王守仁文集》}{《国朝献征录》}
\BookListLine{《唐荆川集》}{《万历三大征考》}
\BookListLine{《土夷考》}{《万历武功录》}
\BookListLine{《宋史》}{《国朝献征录》}
\BookListLine{《新唐书》}{《行边纪闻》}
\BookListLine{}{《炎徼纪闻》}
\par

\noindent 异同各半，四库本《明史·土司传考证》的征引未见其广。

四库本《明史·土司传考证》计153则，中华书局标校本《明史·土司传·校勘记》计146则，两者对勘，相关者仅22则，还不到六分之一。这一方面是由于四库本为「考证」，中华本为「校勘」，两者性质不同，范围自异，因而只有部分的交集；另一方面也因为中华本校勘时未能参考四库本的考证，未能吸纳其成果所致。两者相异的部分正可互补。四库本的考证内容也包括校勘，因此中华本有而四库本无的校勘各条，几乎都可视为四库本的失考；同理，四库本有而中华本无的校勘内容，也可视为中华本的失校。在相关的22则中，涉及年代的有5则、人物人名4则、官名1则、地名7则、隶属5则。总计这相关的22则，两书征引的资料约各为十余种，相同的部分主要为《明实录》、《明一统志》、《明史·地理志》、《汉书·地理志》、《元史·地理志》、《读史方舆纪要》、《明史》相关各传；差异的部分则在四库本用了《土夷考》、《唐荆川集》、《明史纪事本末》、《宋史》，而标校本用了《国榷》、《寰宇通志》、《行边纪闻》。即使考校相同问题，二书考校重点与引用资料也有不同，例如：卷311，建昌卫传裕噜特穆尔之乱诸州皆废革永乐元年复置条，四库本只引《地理志》，中华本增引《太宗实录》；松藩卫传十年黑虎塞贼首多儿太伏诛条，四库本只引《明实录》，中华本增引《国榷》；松藩卫传宣德七年升龙州为宣抚司条，四库本只引《地理志》与《明一统志》，中华本增引《寰宇通志》；天全卫传可自碉门出柘木场条，四库本引《读史方舆纪要》，中华本引《太祖实录》，四库本考证较佳。卷312，永宁卫传樊龙分兵攻合江纳溪条，四库本引《地理志》与


% 原书第 621 页（PDF 第 634 页）
\BookPage{621}{634}
\noindent 《明一统志》，中华本只引《地理志》。卷313，孟定传洪武三十五年土酋刀名扛来朝贡方物赐绮帛钞币设孟定府条，四库本只引《明实录》，中华本增引《地理志》。卷314，丽江传宋时么些蛮蒙醋据之条，四库本只引《明一统志》，中华本增引《寰宇通志》。卷318，结伦州传旧名那兜条，四库本引《明一统志》与《土夷考》，中华本引《读史方舆纪要》；结伦州传宋属结安峒条，四库本引《读史方舆纪要》，中华本引《寰宇通志》、《明一统志》；上石西州传宋属太平条，四库本引《宋史》，中华本引《寰宇通志》、《明一统志》；镇安府传帝命两广总督陶谐治之条，四库本引《陶谐传》、《明实录》、《明史纪事本末》，明显失考，中华本引《明实录》、《田州传》；田州府传斩猛长子于工尧隘条，四库本引《明一统志》、《唐荆川集》，中华本引《归顺州传》、《行边纪闻》；田州传巡按御史曾守约以闻帝命守仁亟为勘处条，四库本照镇安府传帝命两广总督陶谐治之条改，误，中华本引《明实录》；卷319，利州传正统元年泗城岑豹侵据利州地条，四库本引《泗城州传》，中华本引《明实录》；思陵州传以崖州吉阳军儋州南宁军万州万安军俱为州南建州为定安县条，四库本引《明一统志》、《地理志》，考述较详，中华本增引《明实录》。这些情形有些尚待进一步的厘正，有些则可轻易判出高下。

四

王颂蔚《明史考证捃逸·序》云：「汰其文义复沓及空衍无关宏旨者，如云某事某书无考及增某字删某字以符匀刻之类」，我们无法找出王颂蔚当年所看到的进呈本、稿本、正本、初刊样本再作比对，看看他的删汰与择别如何，但是从四库本《明史》的考证文中，仍可观察到一些情况。初步统计，《明史考证捃逸》土司传的部份有280则，四库本《明史》的土司传考证有153则，可见四库本删汰更多，许多考订结果直接在传文中改正，不再于卷末说明，甚至出现将应保留的考证说明也删去的情形。但是，四库本《明史·土司传》的卷末考证文字仍然残留一些王颂蔚所要删汰的内容，卷317平乐传乃定善后六策以闻条「臣章宗瀛按所定善后六策明实录诸书均无考谨识」就是典型的例子。此外，「增减字以符匀刻」的情形不少，馆臣的考证文字也不够精简。


% 原书第 622 页（PDF 第 635 页）
\BookPage{622}{635}
王颂蔚的择别去取也不是全无问题，他的儿子王季烈曾经将《明史考证捃逸》与文津阁四库本《明史》列传卷末的考证逐条互勘，发现二书互见者十之六七，互异者各占十之三四，于是他又在《捃逸》未载而见于文津阁本的明史考证内容选得三十余条录为一卷，作为《捃逸》补遗，其中属于土司传的有四则，并见于文渊阁本，按其内容皆属应保留而不宜删汰者。

四库本《明史·土司传》卷末的考证有153则，王颂蔚《明史考证捃逸》土司传考证有280则，经由比对，四库本153则中有146则见于《捃逸》，另有4则见于《捃逸补遗》，就〈土司传〉的情况而论，《捃逸》的出现几可完全取代四库本《明史》的卷末考证，甚至内容更为丰富，这也许是四库本《明史·土司传考证》不受重视的另一原因。

五

\begin{BookQuote}
《明史·土司传》的内容，聚焦在沿革与叛服二事，这是当时修史者对「土司
\end{BookQuote}

的理解；四库本《明史·土司传》的考证亦不脱这个观念，考证内容首重沿革，次言叛服，征引资料亦以此为标的。但即使依其观念，其所征引的资料范围仍很有限，只用了十余种文献进行核对，其中较常用到的不超过十种；无论如何，谈不上广征博引、考校精当，距当时乾嘉考据学风的水平有一段距离，不宜仅因王颂蔚说有四种稿本，便遽以推断当初馆臣考证过程缜密。个中原因，值得进一步探究，从本文所举〈土司传〉考证内容文字的失误情形而言，馆臣的学术水平与工作态度显然是不能排除的因素。

清代官修《明史》所获评价甚佳，但《明史》的考证从纂修《四库全书》时即已开始，之后代有其人，时有新作，可见问题仍多。时至今日，由于学术观念的进步与学术研究资料的调查整理，土司研究与〈土司传〉的考证早已超越前人的侗限，本文未敢奢言土司研究，只是回顾一段《明史·土司传》的考证经过与内容，谈谈个人的观察与感想而已。
